%------------------------------------------------------------------------------%
% Luis A. D'Afonseca
%
%------------------------------------------------------------------------------%
\documentclass[a4paper,12pt,fleqn]{article}

\usepackage{style_avaliacao}

\ShowAnswers

\newcommand{\D}[2]{\frac{\partial {#1}}{\partial {#2}}}
\newcommand{\DS}[2]{\frac{\partial^2 {#1}}{\partial {#2}^2}}
\newcommand{\DM}[3]{\frac{\partial^2 {#1}}{\partial {#2} \partial {#3}}}

\newcommand{\vetor}[2]{\left(\begin{array}{c} #1 \\ #2 \end{array}\right)}

\newcommand{\veci}{\bm{i}}
\newcommand{\vecj}{\bm{j}}
\newcommand{\veck}{\bm{k}}

\newcommand{\comentario}[1]{{\color{magenta} #1 \par\vspace{\baselineskip}\par}}

%------------------------------------------------------------------------------%
\begin{document}

\makeHeader{CFVVI}{Prova 3 -- Turma 4}{09/09/2024}

%\comentario{
%  Nos exercícios sobre números complexos:
%  \begin{itemize}
%    \item Como esta é uma disciplina de cálculo, exigi que as respostas sejam sempre
%    em radianos, mas os alunos podem fazer os passos intermediários em graus.
%
%    \item No cálculo das raizes o aluno pode deixar a resposta final na forma polar,
%    não é necessário retornar para a forma canônica.
%  \end{itemize}
%}

% 01
%------------------------------------------------------------------------------%
\exercise{50}
Encontre o ponto da esfera \(x^2+y^2+z^2=4\) mais distante do ponto \((1, -1, 1)\).\\
\textit{Dica: não é necessário avaliar completamente a distância para determinar o máximo.}

\begin{answer}
  Queremos minimizar a distância
  \[
    d(x, y, z) = \sqrt{(x-1)^2 + (y+1)^2 + (z-1)^2}
  \]
  Como a função raiz quadrada é crescente podemos minimizar a função
  \[
    f(x, y, z) = (x-1)^2 + (y+1)^2 + (z-1)^2
  \]
  suas derivadas parciais são
  \[
    f_x(x, y, z) = 2(x-1)
    \qquad
    f_y(x, y, z) = 2(y+1)
    \qquad
    f_z(x, y, z) = 2(z-1)
  \]
  As derivadas parciais da função \(g(x, y, z)=x^2+y^2+z^2\) são
  \[
    g_x(x, y, z) = 2x
    \qquad
    g_y(x, y, z) = 2y
    \qquad
    g_z(x, y, z) = 2z
  \]
  Aplicando os Multiplicadores de Lagrange, precisamos resolver o sistema
  \[
  \begin{cases}
    2(x-1) = 2\lambda x \\
    2(y+1) = 2\lambda y \\
    2(z-1) = 2\lambda z \\
    x^2 + y^2 + z^2 = 4
  \end{cases}
  \]
  Das três primeiras equações obtemos
  \begin{align*}
    2(x-1)       & = 2\lambda x          & 2(y+1)       & = 2\lambda y          & 2(z-1)       & = 2\lambda z          \\
    x-1          & = \lambda x           & y+1          & = \lambda y           & z-1          & = \lambda z           \\
    x(1-\lambda) & = 1                   & y(1-\lambda) & = -1                  & z(1-\lambda) & = 1                   \\
    x            & = \frac{1}{1-\lambda} & y            & =\frac{-1}{1-\lambda} & z            & = \frac{1}{1-\lambda}
  \end{align*}
  Sabemos que \(1-\lambda\neq 0\) pois \(x(1-\lambda) = 1\).

  Substituindo na última equação
  \begin{align*}
    x^2 + y^2 + z^2 & = 4
    \\[2mm]
    \left(\frac{1}{1-\lambda}\right)^2 +
    \left(\frac{-1}{1-\lambda}\right)^2 +
    \left(\frac{1}{1-\lambda}\right)^2 & = 4
    \\[2mm]
    \frac{1 + 1 + 1}{\left(1-\lambda\right)^2} & = 4
    \\[2mm]
    \frac{1}{\left(1-\lambda\right)^2} & = \frac{4}{3}
    \\[2mm]
    \frac{1}{1-\lambda} & = \pm\sqrt{\frac{4}{3}}
    = \pm\frac{2}{\sqrt{3}}
  \end{align*}
  Para \(\frac{1}{1-\lambda} = \frac{2}{\sqrt{3}}\)
  \[
  x = \frac{2}{\sqrt{3}}
  \qquad
  y = \frac{-2}{\sqrt{3}}
  \qquad
  z = \frac{2}{\sqrt{3}}
  \]
  Para \(\frac{1}{1-\lambda} = -\frac{2}{\sqrt{3}}\)
  \[
  x = \frac{-2}{\sqrt{3}}
  \qquad
  y = \frac{2}{\sqrt{3}}
  \qquad
  z = \frac{-2}{\sqrt{3}}
  \]
  Temos dois candidatos
  \(\left( \frac{2}{\sqrt{3}}, \frac{-2}{\sqrt{3}}, \frac{2}{\sqrt{3}}\right)\) e
  \(\left( \frac{-2}{\sqrt{3}}, \frac{2}{\sqrt{3}}, \frac{-2}{\sqrt{3}}\right)\).
  Avaliando a função em cada ponto temos
  \begin{align*}
    f\left( \frac{2}{\sqrt{3}}, \frac{-2}{\sqrt{3}}, \frac{2}{\sqrt{3}}\right)
    & = \left(\frac{ 2}{\sqrt{3}}-1\right)^2
    + \left(\frac{-2}{\sqrt{3}}+1\right)^2
    + \left(\frac{ 2}{\sqrt{3}}-1\right)^2 \\[2mm]
    & = \left(\frac{2}{\sqrt{3}}-1\right)^2
    + \left(\frac{2}{\sqrt{3}}-1\right)^2
    + \left(\frac{2}{\sqrt{3}}-1\right)^2 \\[2mm]
    & = 3\left(\frac{2}{\sqrt{3}}-1\right)^2
  \end{align*}
  \begin{align*}
    f\left( \frac{-2}{\sqrt{3}}, \frac{2}{\sqrt{3}}, \frac{-2}{\sqrt{3}}\right)
    & = \left(\frac{-2}{\sqrt{3}}-1\right)^2
    + \left(\frac{ 2}{\sqrt{3}}+1\right)^2
    + \left(\frac{-2}{\sqrt{3}}-1\right)^2 \\[2mm]
    & = \left(\frac{2}{\sqrt{3}}+1\right)^2
    + \left(\frac{2}{\sqrt{3}}+1\right)^2
    + \left(\frac{2}{\sqrt{3}}+1\right)^2 \\[2mm]
    & = 3\left(\frac{2}{\sqrt{3}}+1\right)^2
  \end{align*}
  Portanto o ponto mais distante é
  \(\left( \frac{-2}{\sqrt{3}}, \frac{2}{\sqrt{3}}, \frac{-2}{\sqrt{3}}\right)\).
  \clearpage
\end{answer}

% 02
%------------------------------------------------------------------------------%
\exercise{25}
Dado \(z=\frac{ \left(1-\sqrt{3}\right) + \left(1+\sqrt{3}\right)i}{2+2i} \), calcule
\begin{enumerate}[label=\alph*),itemsep=0pt]
  \item a parte real de $z$,
  \item a parte imaginária de $z$,
  \item o módulo de $z$,
  \item o argumento de $z$
\end{enumerate}

\begin{answer}
  Avaliando $z$
  \begin{align*}
    z
    & = \frac{ \left(1-\sqrt{3}\right) + \left(1+\sqrt{3}\right)i}{2+2i} \\[2mm]
    & = \frac{ \left[\left(1-\sqrt{3}\right) + \left(1+\sqrt{3}\right)i\right](2-2i)}{(2+2i)(2-2i)} \\[2mm]
    & = \frac{ \left(1-\sqrt{3}\right)(2-2i) + \left(1+\sqrt{3}\right)(2i+2)}{2^2 + 2^2} \\[2mm]
    & = \frac{ 2-2i-2\sqrt{3}+2\sqrt{3}i + 2i+2+2\sqrt{3}i +2\sqrt{3}}{8} \\[2mm]
    & = \frac{ 4 + 4\sqrt{3}i}{8} \\[2mm]
    & = \frac{1}{2} + \frac{\sqrt{3}}{2}i
  \end{align*}
  Parte real de $z$
  \[
    \Re(z) = \frac{1}{2}
  \]
  Parte imaginária de $z$
  \[
    \Im(z) = \frac{\sqrt{3}}{2}
  \]
  Módulo de $z$
  \[
    \abs{z}
    = \sqrt{\left(\frac{1}{2}\right)^2 + \left(\frac{\sqrt{3}}{2}\right)^2}
    = \sqrt{\frac{1}{4} + \frac{3}{4}}
    = \sqrt{\frac{1+3}{4}}
    = \sqrt{1}
    = 1
  \]
  Argumento de $z$, \(\varphi=\arg(z)\)
  \[
    \sin(\varphi)
    = \frac{\Im(z)}{\abs{z}}
    = \frac{\sqrt{3}}{2}
  \qquad\qquad
    \cos(\varphi)
    = \frac{\Re(z)}{\abs{z}}
    = \frac{1}{2}
  \]
  \[
    \arg(z)
    = \ang{60}
    = \frac{\pi}{3}
  \]
  \clearpage
\end{answer}

% 03
%------------------------------------------------------------------------------%
\exercise{25}
Encontre as raízes cúbicas complexas do número \( z = 1 + i \).

\begin{answer}
  Escrevemos \(z = 1 + i\) na forma trigonométrica. O módulo é
  \[
    \rho = \abs{1+i} = \sqrt{1^2 + 1^2} = \sqrt{2}
  \]
  O argumento é
  \[
    \varphi = \arg(1+i) = \ang{45} = \frac{\pi}{4}
  \]
  \[
    z = \sqrt{2} \left[ \cos\left(\frac{\pi}{4}\right) + i \sin\left(\frac{\pi}{4}\right) \right]
  \]
  As raízes cúbicas são
  \begin{align*}
    u_k
    & = \sqrt[3]{\rho} \left[
        \cos\left(\frac{\varphi + 2k\pi}{3}\right)
        + i \sin\left(\frac{\varphi + 2k\pi}{3}\right)
    \right]
    & k = 0, 1, 2    \\[2mm]
    & = \sqrt[3]{\sqrt{2}} \left[
        \cos\left(\frac{1}{3}\left(\frac{\pi}{4} + 2k\pi\right)\right)
        + i \sin\left(\frac{1}{3}\left(\frac{\pi}{4} + 2k\pi\right)\right)
        \right] \\[2mm]
    & = \sqrt[6]{2} \left[
        \cos\left( \frac{\pi + 8k\pi}{12} \right)
        + i \sin\left( \frac{\pi + 8k\pi}{12} \right)
        \right]
  \end{align*}
  Para cada valor de $k$
  \begin{align*}
    u_0
    & = \sqrt[6]{2} \left[
        \cos\left( \frac{\pi + 8\times0\pi}{12} \right)
        + i \sin\left( \frac{\pi + 8\times0\pi}{12} \right)
        \right] \\[2mm]
    & = \sqrt[6]{2} \left[
        \cos\left( \frac{\pi}{12} \right)
        + i \sin\left( \frac{\pi}{12} \right)
        \right]
  \end{align*}
  \begin{align*}
    u_1
    & = \sqrt[6]{2} \left[
        \cos\left( \frac{\pi + 8\times1\pi}{12} \right)
        + i \sin\left( \frac{\pi + 8\times1\pi}{12} \right)
        \right] \\[2mm]
    & = \sqrt[6]{2} \left[
        \cos\left( \frac{9\pi}{12} \right)
        + i \sin\left( \frac{9\pi}{12} \right)
        \right] \\[2mm]
    & = \sqrt[6]{2} \left[
        \cos\left( \frac{3\pi}{4} \right)
        + i \sin\left( \frac{3\pi}{4} \right)
        \right] \\[2mm]
    & = \sqrt[6]{2} \left[
        - \frac{\sqrt{2}}{2}
        + i \frac{\sqrt{2}}{2}
        \right] \\[2mm]
    & = - \frac{2^{\sfrac{1}{6}}}{2^{\sfrac{1}{2}}}
        + \frac{2^{\sfrac{1}{6}}}{2^{\sfrac{1}{2}}} i  \\[2mm]
    & = - 2^{-\sfrac{1}{3}}
        + 2^{-\sfrac{1}{3}} i \\[2mm]
    & = - \frac{1}{\sqrt[3]{2}}
        + \frac{i}{\sqrt[3]{2}}
  \end{align*}
  \begin{align*}
    u_2
    & = \sqrt[6]{2} \left[
        \cos\left( \frac{\pi + 8\times2\pi}{12} \right)
        + i \sin\left( \frac{\pi + 8\times2\pi}{12} \right)
        \right] \\[2mm]
    & = \sqrt[6]{2} \left[
        \cos\left( \frac{17}{12}\pi \right)
        + i \sin\left( \frac{17}{12}\pi \right)
        \right]
  \end{align*}
\end{answer}

%------------------------------------------------------------------------------%
\end{document}
%------------------------------------------------------------------------------%
